\documentclass[12pt]{article}

\usepackage[utf8]{inputenc}
\usepackage[T1]{fontenc}
\usepackage[polish]{babel}
\usepackage{listings}
\usepackage{xcolor}
\usepackage{geometry}
\geometry{margin=2.5cm}

\lstset{
    language=C,
    basicstyle=\ttfamily\small,
    keywordstyle=\color{blue},
    stringstyle=\color{red},
    commentstyle=\color{green!40!black},
    numbers=left,
    numberstyle=\tiny,
    frame=single,
    breaklines=true
}

\title{Dokumentacja końcowa projektu\\planar-graph-visualization}
\author{Mateusz Kamieniecki \\ Jan Rękas}
\date{\today}

\begin{document}

\maketitle
\tableofcontents

\section{Co zostało osiągnięte w projekcie}
W projekcie dostarczono działającą aplikację CLI w języku C, która realizuje pełny
pipeline:
\begin{center}
	wejście grafu \textrightarrow{} obliczenie layoutu \textrightarrow{} zapis wyników
\end{center}
Zaimplementowano i zintegrowano:
\begin{itemize}
	\item wczytanie grafu z pliku tekstowego,
	\item algorytm \texttt{1} (Tutte),
	\item algorytm \texttt{2} (Spectral layout),
	\item zapis wyników do formatu binarnego i tekstowego,
\end{itemize}

\section{Uzyskane cechy funkcjonalne}
\subsection{Wejście}
Każda linia wejścia opisuje jedną krawędź:
\begin{verbatim}
<nazwa_krawędzi> <from> <to> <waga>
\end{verbatim}
W projekcie przyjęto i wdrożono 1-indeksowanie identyfikatorów wierzchołków
(\texttt{1..N}); wierzchołek \texttt{0} nie jest używany.

\subsection{Wyjście}
\subsubsection{Format tekstowy (\texttt{.txt})}
\begin{verbatim}
<liczba_wierzchołków>
<id> <x> <y>
...
\end{verbatim}

\subsubsection{Format binarny (\texttt{.bin})}
Kolejno zapisywane są:
\begin{itemize}
	\item liczba wierzchołków: \texttt{int} (4 bajty),
	\item dla każdego wierzchołka: \texttt{id (int)}, \texttt{x (double)}, \texttt{y (double)}.
\end{itemize}

\subsection{Interfejs CLI}
Program:
\begin{verbatim}
./bin/planar_graph [flagi]
\end{verbatim}

\begin{center}
	\begin{tabular}{|c|c|p{7.5cm}|}
		\hline
		\textbf{Flaga} & \textbf{Argument}  & \textbf{Opis}                                                        \\
		\hline
		\texttt{-i}    & \texttt{<ścieżka>} & Plik wejściowy (domyślnie \texttt{input.txt})                        \\
		\hline
		\texttt{-o}    & \texttt{<katalog>} & Katalog wyjściowy (domyślnie bieżący)                                \\
		\hline
		\texttt{-n}    & \texttt{<nazwa>}   & Nazwa pliku wyjściowego bez rozszerzenia (domyślnie \texttt{output}) \\
		\hline
		\texttt{-a}    & \texttt{1|2}       & Wybór algorytmu (\texttt{1} = Tutte, \texttt{2} = Spectral Layout)   \\
		\hline
		\texttt{-b}    & —                  & Zapis binarny                                                        \\
		\hline
		\texttt{-h}    & —                  & Zapis tekstowy                                                       \\
		\hline
	\end{tabular}
\end{center}

W ramach finalizacji osiągnięto:
\begin{itemize}
	\item domyślny zapis binarny, gdy nie podano \texttt{-b} ani \texttt{-h},
	\item możliwość zapisu obu formatów jednocześnie przy \texttt{-b -h}.
\end{itemize}

\section{Uzyskany stan implementacji}
\subsection{Architektura modułów}
\begin{itemize}
	\item \texttt{main.c} — parsowanie flag, walidacja wejścia, uruchomienie algorytmu, zapis wyników.
	\item \texttt{fileio.c/.h} — odczyt grafu i zapis layoutu.
	\item \texttt{graph.c/.h} — struktury grafu i operacje na listach sąsiedztwa.
	\item \texttt{tutte.c/.h} — implementacja osadzania Tutte.
	\item \texttt{node.h} — definicje \texttt{Node} i \texttt{GraphLayout}.
	\item \texttt{utils.c/.h} — parsowanie wyboru algorytmu.
\end{itemize}

\subsection{Kluczowe struktury danych}
\begin{lstlisting}
typedef struct Edge {
    int to;
    double weight;
    char *name;
} Edge;

typedef struct Vector {
    Edge *data;
    int size;
    int capacity;
} Vector;

typedef struct Graph {
    int vertices_num;
    Vector *adj;
} Graph;

typedef struct Node {
    int id;
    double x;
    double y;
} Node;

typedef struct {
    int count;
    Node *nodes;
} GraphLayout;
\end{lstlisting}

\subsection{Najważniejsze interfejsy}
\subsubsection{\texttt{find\_max\_vertex\_id}}
\textbf{Przyjmuje:} \texttt{const char *filepath}\\
\textbf{Zwraca:} \texttt{int} (\texttt{-1} przy braku dostępu do pliku, \texttt{0} dla pustych/niepoprawnych danych, inaczej maksymalne ID)

\subsubsection{\texttt{load\_graph\_from\_file}}
\textbf{Przyjmuje:} \texttt{const char *filepath}, \texttt{int vertices\_amount}\\
\textbf{Zwraca:} \texttt{Graph*}

\subsubsection{\texttt{add\_edge}}
\textbf{Przyjmuje:} \texttt{Graph *graph}, \texttt{int from}, \texttt{int to}, \texttt{double weight}, \texttt{char *name}\\
\textbf{Zwraca:} \texttt{void}

\subsubsection{\texttt{find\_embedding}}
\textbf{Przyjmuje:} \texttt{Graph *g}, \texttt{GraphLayout *layout}\\
\textbf{Zwraca:} \texttt{void}

\section{Zrealizowana obsługa błędów i ostrzeżeń}
W kodzie i dokumentacji ujednolicono komunikaty. Obsłużone są m.in.:
\begin{itemize}
	\item nieznane/niekompletne flagi (\texttt{EXIT\_FAILURE}),
	\item nieobsługiwany algorytm (\texttt{EXIT\_FAILURE}),
	\item brak dostępu do pliku wejściowego (\texttt{EXIT\_FAILURE}),
	\item brak dodatnich identyfikatorów wierzchołków (\texttt{EXIT\_FAILURE}),
	\item ostrzeżenia o pomijanych krawędziach (program działa dalej),
	\item problemy alokacji pamięci (zakończenie awaryjne),
	\item brak cyklu / problem numeryczny w Tutte (komunikat informacyjny).
\end{itemize}

\section{Kompilacja i uruchomienie programu}
Program kompilowany jest przez \texttt{Makefile} znajdujący się w katalogu głównym repozytorium.

\subsection{Kompilacja}
\begin{verbatim}
make build
\end{verbatim}
Po kompilacji plik wykonywalny znajduje się pod ścieżką:
\begin{verbatim}
./bin/planar_graph
\end{verbatim}

\subsection{Uruchomienie}
Uruchomienie domyślne (z domyślnym \texttt{input.txt}):
\begin{verbatim}
./bin/planar_graph
\end{verbatim}

Uruchomienie z przykładowymi flagami:
\begin{verbatim}
./bin/planar_graph -i graf.txt -o wyniki -n wynik -a 1 -b -h
\end{verbatim}

\section{Limitacje}
Obecna wersja programu ma następujące ograniczenia:
\begin{itemize}
	\item algorytm Tutte może nie zadziałać dla danego grafu z powody jego konstrukcji (brak cyklu, nie 3-connected)
\end{itemize}

\section{Podsumowanie końcowe}
Finalnie uzyskano spójny, działający program CLI do generowania layoutu grafów,
z jednolitym 1-indeksowaniem wierzchołków, poprawnym zapisem wyników
(\texttt{.txt} i \texttt{.bin}), zaktualizowaną dokumentacją oraz
komunikatami błędów/ostrzeżeń
\end{document}
